\subsection{Cơ chế hoạt động}
\label{sec:2-mechanism}

\subsubsection{Kiến trúc pipeline}
\label{sec:2-pipeline}

Một pipeline micro-batch điển hình cho dữ liệu cảm biến gồm các thành phần nối tiếp trong Hình~\ref{fig:2-pipeline}. Dữ liệu từ cảm biến qua cổng kết nối (gateway) vào nhật ký sự kiện Kafka. Kafka tiếp tục nhận sự kiện trong khi Spark đang tính toán, nhờ đó khâu thu thập và khâu xử lý tách rời nhau. Kết quả được ghi vào cơ sở dữ liệu phục vụ bảng điều khiển. Tiến độ xử lý và trạng thái được lưu bền vững trong checkpoint.

\begin{figure}[H]
\centering
\begin{tikzpicture}[
  node distance=0.55cm,
  box/.style={rectangle, rounded corners, draw, minimum width=2.0cm, minimum height=1.0cm, align=center, font=\small, fill=gray!10},
  arr/.style={-{Stealth}, thick}]
\node[box] (s) {Cảm biến};
\node[box, right=of s] (g) {Gateway};
\node[box, right=of g, fill=blue!10] (k) {Kafka\\nhật ký sự kiện};
\node[box, right=of k, fill=orange!15] (p) {Spark\\tính toán};
\node[box, right=of p] (d) {CSDL \&\\dashboard};
\node[box, below=0.9cm of p, fill=green!10] (c) {Checkpoint\\offset + trạng thái};
\draw[arr] (s)--(g); \draw[arr] (g)--(k); \draw[arr] (k)--(p); \draw[arr] (p)--(d);
\draw[arr, dashed] (p)--(c);
\end{tikzpicture}
\caption{Kiến trúc pipeline micro-batch từ cảm biến đến bảng điều khiển}
\label{fig:2-pipeline}
\end{figure}

\subsubsection{Vòng đời xử lý của một micro-batch}
\label{sec:2-lifecycle}

Mỗi micro-batch trải qua sáu bước:
\begin{enumerate}
  \item \textbf{Trigger}: bắt đầu lượt chạy tiếp theo.
  \item \textbf{Lập kế hoạch và ghi offset}: chọn các bản ghi mới và ghi lại phạm vi offset sẽ đọc.
  \item \textbf{Đọc và biến đổi}: áp dụng truy vấn lên dữ liệu của lô.
  \item \textbf{Cập nhật trạng thái}: dành cho các truy vấn có trạng thái (cửa sổ, tổng hợp, khử trùng lặp).
  \item \textbf{Ghi kết quả}: gửi đầu ra tới nơi lưu (sink).
  \item \textbf{Cam kết lô}: ghi nhận lô đã hoàn tất.
\end{enumerate}

\subsubsection{Chu kỳ trigger và ví dụ dữ liệu nông nghiệp}
\label{sec:2-trigger}

Chu kỳ trigger quyết định độ chi tiết của lô. Giả sử một cảm biến độ ẩm đất gửi số đo lên Kafka và trigger chạy mỗi 5 giây (Bảng~\ref{tab:2-example}). Mỗi lô tính một giá trị tổng hợp rồi lưu kết quả, hiển thị trên bảng điều khiển và đưa vào các quy tắc cảnh báo.

\begin{table}[H]
\centering
\caption{Ví dụ hai micro-batch 5 giây của một cảm biến độ ẩm đất}
\label{tab:2-example}
\small
\begin{tblr}{
  colspec = {X[0.4,l]X[1,l]X[2.2,l]X[1,l]}, rows = {abovesep=2pt, belowsep=2pt},
  hlines, vlines, row{1} = {font=\bfseries, bg=gray!20},
}
Lô & Trigger & Số đo đến Kafka (độ ẩm \%) & Trung bình \\
1 & 08:00:05 & 22 (08:00:01), 24 (08:00:03), 26 (08:00:04) & 24\% \\
2 & 08:00:10 & 28 (08:00:07), 30 (08:00:09) & 29\% \\
\end{tblr}
\end{table}

\subsubsection{Độ tin cậy: phục hồi sau lỗi và dữ liệu đến muộn}
\label{sec:2-reliability}

\paragraph{Phục hồi sau lỗi.}
Checkpoint lưu lại một cách bền vững tiến độ xử lý, gồm các offset đã đọc và trạng thái của truy vấn. Khi khởi động lại, hệ thống phát lại và hoàn tất lô chưa được cam kết nên không làm mất hay bỏ sót dữ liệu.

\paragraph{Dữ liệu đến muộn.}
Cần phân biệt thời gian sự kiện (event time) với thời gian đến (arrival time). Một số đo có thể được tạo lúc 08:00:03 nhưng đến hệ thống sau đó vì mất kết nối. Để nhận được các số đo đến muộn này, hệ thống phải giữ lại trạng thái của các cửa sổ chưa đóng. \emph{Watermark} là mốc đánh dấu tiến độ của thời gian sự kiện, không phải bộ đếm giờ chạy theo đồng hồ thực. Mỗi khi watermark được đẩy sang mốc thời gian mới hơn, hệ thống dọn trạng thái của những cửa sổ đã quá thời hạn chờ số đo muộn, nhờ đó bộ nhớ luôn ở mức giới hạn.

\paragraph{Điều kiện để đạt exactly-once.}
``Đúng một lần'' (exactly-once) nghĩa là dù có sự cố và phải xử lý lại, mỗi sự kiện vẫn chỉ ảnh hưởng đến kết quả đúng một lần. Muốn được như vậy quá trình xử lý cần phải thoả mãn đồng thời ba điều kiện: nguồn có thể phát lại (replayable) để đọc lại đúng các sự kiện của lô bị gián đoạn; engine có checkpoint để biết lô nào đã hoàn tất; nơi ghi chịu được việc ghi lặp, hoặc nhờ thao tác idempotent (ghi lại nhiều lần vẫn cho cùng kết quả), hoặc nhờ giao dịch (transactional) để một lô chỉ có hiệu lực khi được cam kết trọn vẹn. Nếu thiếu điều kiện thứ ba, chẳng hạn khi ghi vào một sink Kafka thông thường không dùng giao dịch, hệ thống chỉ đạt at-least-once: không mất sự kiện nhưng một sự kiện có thể bị xử lý nhiều hơn một lần.
